\section{مدیریت کارهای مسدود شده (\lr{Blocked Items}) و گلوگاه‌ها}
یک اصل کلیدی در رویکرد \lr{Lean DAD} این است که محدودیت‌های \lr{WIP} باید کارهای مسدود شده را نیز شامل شوند:
\begin{itemize}
	\item کارت‌هایی که به دلیل وابستگی به ماژول دیگر یا نیاز به یک \lr{Spike} فنی مسدود (\lr{Blocked}) می‌شوند، \textbf{همچنان بخشی از ظرفیت \lr{WIP} آن ستون محسوب می‌شوند.}
	\item این سیاست تیم را مجبور می‌کند که به جای رها کردن کار مسدود شده و کشیدن کارت جدید (که باعث نقض \lr{WIP} می‌شود)، با رویکرد \lr{Swarming} روی رفع مانع تمرکز کند.
	\item کارت‌های مسدود شده با برچسب قرمز مشخص شده و در جلسات هماهنگی روزانه (\lr{Daily Coordination})، بالاترین اولویت بررسی را دارند.
\end{itemize}